\chapter{Conclusion and Future Work}
\label{ch:conclusion}

\section{Conclusion}
\label{sec:conclusion-summary}

This thesis studied multimodal sensing-assisted user association and tracking
for beam management in vehicular networks. The motivating problem is that
vehicle users in road-side communication scenarios experience rapidly changing
geometry, blockage, and beam quality, while repeated beam sweeping and channel
measurement can consume non-negligible radio resources. To address this
setting, the thesis considered a BS-centric sensing infrastructure in which the
serving base station (BS) and nearby sensing nodes (SNs) provide multiview
camera observations, and the BS also provides an ISAC point cloud. These
observations are organised through a bird's-eye-view (BEV) representation so
that vehicle detection, per-target beam prediction, and communication-oriented
evaluation can be studied within a common spatial frame.

The first part of the work formulated a multimodal distributed
sensing-assisted beam management framework. In this framework, camera modality
and ISAC modality are fused in BEV space to detect vehicle targets and estimate
per-target beam-class posteriors. The resulting ranked Top-K candidate beams
provide a sensing-driven interface between the perception system and the subsequent
communication evaluation. The objective of this design is not to replace
communication measurements entirely, but to use sensing observations to reduce
the frequency and scope of expensive beam measurements when the sensing output
is sufficiently reliable.

The second part of the work addressed the identity mismatch between sensing
targets and communication users. Since a detected vehicle target is not
automatically associated with a communication user identity, the proposed
workflow uses the initial SSB beam sweep power spectrum to establish
target-user association. After this initial binding, an interacting multiple
model (IMM) tracking process maintains the temporal continuity of vehicle
states and supports beam maintenance in subsequent frames. This association and
tracking layer is central to the thesis scope because it connects per-frame
sensing predictions with a persistent communication service process.

The third part of the work specified a multimodal simulation and system-level
evaluation methodology. The evaluation connects road traffic
simulation, multiview sensing data generation, ISAC-like point clouds,
ray-tracing-based beam labels, and simplified communication metrics. Within
this abstraction, beam prediction accuracy, target-user association, tracking
quality, beam training overhead, beam mismatch, zero-forcing (ZF) precoding
SINR, spectral efficiency, and overhead-adjusted throughput can be considered
together. This evaluation pipeline makes explicit how errors in sensing and
identity maintenance may propagate into communication-level performance.
The overall contribution is therefore a framework-level and evaluation-level
study of multimodal sensing-assisted beam management rather than a
protocol-complete 5G NR implementation.

\section{Limitations}
\label{sec:conclusion-limitations}

Several limitations follow from the current scope of the thesis. First, the
evaluation relies on a simulated or generated multimodal
environment. Such a setting is necessary for controlled access to camera
observations, ISAC-like point clouds, vehicle trajectories, beam labels, and
channel data, but it may introduce a domain gap relative to real road-side
deployments. The appearance of vehicles, sensing noise, calibration quality,
traffic behaviour, occlusion patterns, and propagation conditions in real
networks may differ from those represented by the simulation pipeline.

Second, the communication evaluation is a simplified and interpretable
abstraction rather than a complete 5G NR protocol stack. The thesis models
beam management, beam mismatch, ZF precoding SINR, spectral efficiency, and
overhead-adjusted throughput at a system level, but it does not fully model
protocol-level details such as SSB and CSI-RS scheduling, feedback latency,
control signaling, HARQ behaviour, scheduler interaction, or mobility management
procedures. As a result, throughput-related quantities should be interpreted as
evaluation indicators under the stated abstraction, not as direct predictions
of deployment-level 5G NR performance.

Third, the proposed sensing-assisted workflow depends on the reliability of
multimodal sensing and cross-modal alignment. Camera calibration errors,
time-synchronization errors, sparse or noisy ISAC point clouds, missing
modalities, and severe occlusion can degrade BEV detection, beam prediction,
target-user association, and tracking. The framework includes partial beam
measurement and re-association as possible recovery mechanisms, but the precise
robustness limits and failure thresholds remain important factors for
system-level evaluation and deployment-oriented future work.

Fourth, the thesis focuses primarily on a BS-centric local service region.
This choice keeps the problem tractable and allows the relationship between
multimodal sensing, target-user association, tracking, and beam management to
be studied clearly. However, it leaves broader network-level questions only
partially addressed, including cooperative multi-BS beam management,
inter-cell interference, load-aware service selection, and handover decisions
for vehicles moving across BS coverage regions.

\section{Future Work}
\label{sec:conclusion-future-work}

Future work should extend the thesis in directions that directly correspond to
the limitations above.

\subsection{Real-World Deployment and Domain Gap}
\label{subsec:future-domain-gap}

A natural next step is to evaluate the framework with real road-side sensing
and communication data. This may include collecting synchronized multiview
camera observations, ISAC or radar-like point clouds, vehicle trajectories, and
communication beam measurements from instrumented road-side infrastructure.
Such data would make it possible to quantify the domain gap between simulation
and deployment, validate the assumptions used in the BEV fusion pipeline, and
study whether synthetic-to-real adaptation is necessary for reliable beam
prediction and target-user association.

\subsection{Protocol-Level 5G NR Simulation}
\label{subsec:future-5g-nr}

The communication abstraction can be extended toward a fuller 5G NR
protocol-level simulation. Future work may incorporate more detailed modelling
of SSB and CSI-RS timing, beam reporting, feedback delay, scheduling, control
overhead, retransmission behaviour, and mobility procedures. This would allow
sensing-assisted beam management to be evaluated not only through simplified
SINR and throughput indicators, but also through protocol-aware latency,
resource utilization, and service continuity metrics.

\subsection{Online and Domain Adaptation}
\label{subsec:future-online-adaptation}

The present framework is primarily described as an offline trained and
evaluated system. Future work can investigate online adaptation so that the
beam prediction, association, and tracking modules can adjust to new roads,
sensor placements, vehicle distributions, and propagation conditions. Possible
directions include continual calibration, self-supervised adaptation from
communication measurements, and domain adaptation between simulated and real
scenarios. These extensions should be designed carefully to avoid reinforcing
incorrect associations or unreliable beam predictions during deployment.

\subsection{Cooperative Multi-BS Beam Management and Handover}
\label{subsec:future-multi-bs-handover}

Finally, the BS-centric design can be extended to cooperative multi-BS beam
management and handover. In a multi-BS network, sensing and tracking states may
be shared across neighboring BSs to support service-BS selection, handover
preparation, and coordinated beam decisions for fast-moving vehicles. This
extension would require new association variables linking vehicles, sensing
tracks, serving BSs, and candidate beams, as well as evaluation metrics that
capture handover reliability, inter-cell coordination overhead, and service
continuity.
